% Provenance (never printed): sources/book2/BRIEFS.md, g10-01. Data and models; statements of fact
% and of value; correlation and cause. Every number about a real country comes from
% figdata/grade-10/01-economic-models/ (Maddison Project Database 2023, WDI); Brenland is invented.
\chapter{Economic Models}\label{ch:g10:economic-models}

Imagine a clerk in 1929 with a ruler, a sheet of paper ruled on a times-ten scale and sixty
years of figures. Asked to guess the income of an average American and of an average
Argentine in 2022, ninety-three years later, the clerk does the simplest thing possible: plot
income per person year by year since 1870, lay the ruler along the points, and extend the line.
In 1929 that was not a silly bet. Argentina was then one of the richer countries of the world,
with an income per person of about 58\,\% of the American one, and over the sixty years before,
it had grown a little faster.
% figdata: trend_points (ARG 6961 / USA 11954 = 0.582, 1929); trend_fit (rates 1.64 and 1.53 % a year)

The opening chart shows what happened to the clerk's two lines. For the United States, the value
observed in 2022 was about 17\,\% above the line: on the scale of the chart, close. For
Argentina, what was observed in 2022 was a little more than half of the line's value. The same
method, the same years, two very different results. This chapter is about the method: what a
model of the economy is, how it is stated, how it is tested against what was then observed, and
where it breaks. It ends with the question that the opening chart raises and does not answer:
when two things move together, how can we tell whether one causes the other?
% figdata: trend_fit (USA pct 117: observed / line = 1.17; ARG pct 56)

\begin{openingchart}{Income per person in the United States and Argentina, 1870--2022 (solid),
and the constant-growth line fitted to each country's points of 1870--1929 and then extended to
2022 (dashed), on a times-ten scale. The line is the straight line that passes closest to the
points of 1870--1929; nothing after 1929 was used to draw it. Maddison has no estimates for
Argentina in 1871--1874.}{Maddison Project Database 2023; GDP per person in 2011 dollars at purchasing-power parity}
\begin{tikzpicture}
\begin{semilogyaxis}[omtimesten, width=0.94\linewidth, height=7.4cm, xmin=1865, xmax=2028,
  ymin=1000, ymax=100000, xlabel={year}, ylabel={dollars per person (2011 prices)},
  xtick={1870,1890,...,2010}, legend pos=north west]
\addplot[omDef, thick] table[x=year, y=gdppc] {figdata/grade-10/01-economic-models/out/obs_usa.csv};
\addlegendentry{United States, observed}
\addplot[omThm, thick] table[x=year, y=gdppc] {figdata/grade-10/01-economic-models/out/obs_arg.csv};
\addlegendentry{Argentina, observed}
\addplot[omDef, thick, dashed] table[x=year, y=usa] {figdata/grade-10/01-economic-models/out/trend_lines.csv};
\addlegendentry{United States, line of 1870--1929}
\addplot[omThm, thick, dashed] table[x=year, y=arg] {figdata/grade-10/01-economic-models/out/trend_lines.csv};
\addlegendentry{Argentina, line of 1870--1929}
\draw[black!55, line width=0.6pt] (axis cs:1929.5,1000) -- (axis cs:1929.5,100000);
\node[font=\footnotesize, anchor=south east, text=black!70] at (axis cs:1928,1150) {fitted on these years};
\node[font=\footnotesize, anchor=south west, text=black!70] at (axis cs:1931,1150) {judged on these years};
\end{semilogyaxis}
\end{tikzpicture}
\end{openingchart}

\begin{recall}
From the grades 6--9 volume: the \emph{growth rate} of a quantity is its change over a year as a
percentage of its starting value; growth at the same rate year after year is \emph{compound
growth}, which multiplies the quantity by the same factor $1+g$ every year. On a
\emph{times-ten scale}, each step of the axis multiplies by ten, so compound growth at a fixed rate
draws a straight line. That volume also boxed small models in four parts, and said that a
\emph{correlation} across countries is not yet a cause.
\end{recall}

\section{What a model is}\label{sec:g10:economic-models:what}

A street map leaves out almost everything about a city: the colour of the houses, the people,
the weather, the height of the trees. That is why it is useful. A map drawn at full size, with
every detail, would be as large as the city and of no help in finding the station. A model of the
economy does the same: it keeps a few quantities, states how they are linked, and leaves the rest
out on purpose. What it leaves out is a choice, and the choice is stated.

\begin{definition}[Economic model]\label{def:g10:economic-models:model}
An \emph{economic model}\index{economic model} is a simplified description of part of the economy,
built on stated assumptions, from which predictions about quantities that can be observed are
derived.
\end{definition}

The clerk's ruler is a model. Write $y_t$ for income per person $t$ years after a starting year,
$y_0$ for its value in the starting year, and $g$ for the growth rate, written as a decimal ($g =
0.02$ for 2\,\% a year). The assumption is that $g$ is the same every year. Then each year
multiplies income by $1+g$, and after $t$ years
\begin{equation}\label{eq:g10:economic-models:growth}
  y_t = y_0\,(1+g)^t .
\end{equation}
The model takes $y_0$ and $g$ as given: it does not say where they come from. It explains $y_t$,
for every $t$. Economists name the two roles.

\begin{definition}[Endogenous and exogenous variables]\label{def:g10:economic-models:endogenous}
An \emph{endogenous variable}\index{endogenous variable} is a quantity whose value a model
explains: the model's equations give it once the other quantities are known. An \emph{exogenous
variable}\index{exogenous variable} is a quantity that the model takes as given from outside: the
model uses its value but does not explain it.
\end{definition}
\phantomsection\label{def:g10:economic-models:exogenous}

In \eqref{eq:g10:economic-models:growth}, $y_t$ is endogenous; $y_0$, $g$ and the number of years
$t$ are exogenous. The same quantity can play either role in different models. A model of a
country's saving and investment, met in grade 12, will try to explain the growth rate $g$ that
this one takes as given. No model explains everything; each one draws a line between what it
explains and what it takes from outside, and a good model draws the line where the outside
quantities are either known or can be measured.

\begin{example}[A bakery in Brenland]\label{ex:g10:economic-models:bakery}
Brenland is an invented economy that this series uses for exercises; its currency is the dalen. A
bakery in Brenland sells $q$ loaves a day at a price of $p$ dalens each, and its takings are $R =
p\,q$ dalens a day. In this tiny model $R$ is endogenous, $p$ and $q$ are exogenous. A larger model
of the same bakery might add that fewer loaves are sold when the price is higher; then $q$ becomes
endogenous, explained by $p$, and only $p$ (and whatever else moves the number of buyers) is
exogenous.
\end{example}

A model also needs a clause that ordinary speech leaves implicit. ``A higher price of bread means
fewer loaves sold'' is not true if, in the same week, a heat wave closes the bakery's rival or a
holiday brings visitors to town: the number of loaves sold changes for several reasons at once.
The statement means that the price alone has changed.

\begin{definition}[Other things equal]\label{def:g10:economic-models:other-things-equal}
A prediction holds \emph{other things equal}\index{other things equal} when it describes the effect
of a change in one \omterm{def:g10:economic-models:endogenous}{exogenous variable} while every other \omterm{def:g10:economic-models:endogenous}{exogenous variable} of the model keeps its
value.
\end{definition}

``\omterm{def:g10:economic-models:other-things-equal}{Other things equal}'' is not a way to escape a test. It tells the tester what to look for: a case
where the one variable changed and the others did not, or a way to take the change in the others
into account. Much of this book, and all of the econometrics of the university volumes, is about
finding such cases.

\section{The four-point frame}\label{sec:g10:economic-models:frame}

Every model in this book is summed up in a box with four fields, always in the same order. The
grades 6--9 volume used the same boxes; this section says what each field must contain.

\begin{itemize}
\item \textbf{Assumes.} The assumptions, stated so that one can check whether a real case meets
  them. ``The growth rate is the same every year'' can be checked; ``the economy behaves normally''
  cannot.
\item \textbf{Predicts.} An observable quantity, a direction and, whenever possible, a size. ``Income
  per person will be higher'' is a weak prediction; ``income per person will be about
  \usd{50000} in 2022'' is a strong one, because many observations would contradict it.
\item \textbf{Tested.} Where the prediction was compared with data that were \emph{not} used to
  build the model. A line fitted to the years 1870--1929 always passes close to the points of
  1870--1929: that is how it was drawn, and it proves nothing. The test is what the line says about
  1930--2022, before anyone has seen those years.
\item \textbf{Breaks.} The real cases in which an assumption fails and the prediction with it. A
  model that is never said to break has not been tested.
\end{itemize}

\begin{omfigure}
\begin{tikzpicture}[econ, node distance=4mm and 9mm,
  box/.style={draw=omModel!75!black, fill=omModel!6, rounded corners=1.5mm, minimum width=31mm,
    minimum height=15mm, align=center, text width=29mm, font=\small}]
\node[box] (a) {\textbf{Assumes}\\[1pt] what the model takes as true};
\node[box, right=of a] (p) {\textbf{Predicts}\\[1pt] a quantity, a direction, a size};
\node[box, below=of p] (t) {\textbf{Tested}\\[1pt] against data not used to build it};
\node[box, below=of a] (b) {\textbf{Breaks}\\[1pt] where an assumption fails};
\draw[->, thick, omModel!75!black] (a) -- (p);
\draw[->, thick, omModel!75!black] (p) -- (t);
\draw[->, thick, omModel!75!black] (t) -- (b);
\draw[->, thick, omModel!75!black, dashed] (b) -- node[left, font=\footnotesize, align=right]
  {revise the\\assumptions} (a);
\end{tikzpicture}

{\small The four-point frame. A model is stated by its assumptions, judged by its predictions on data
it was not built from, and known by where it fails; a failure sends the modeller back to the
assumptions.}
\end{omfigure}

To fit the clerk's line by hand, two points are enough: the first and the last year of the fitting
period. If income goes from $y_0$ to $y_T$ in $T$ years at a constant rate, then $y_T = y_0(1+g)^T$,
so $(1+g)^T = y_T/y_0$, and $1+g$ is the number that, multiplied by itself $T$ times, gives the
ratio $y_T/y_0$. A calculator finds it as the power $(y_T/y_0)^{1/T}$.

\begin{method}[A constant-growth forecast from two points]\label{met:g10:economic-models:two-points}
To fit the constant-growth model to the values $y_0$ and $y_T$, $T$ years apart, and project it $s$
years beyond the last one:
\begin{enumerate}
\item compute the ratio $y_T / y_0$;
\item compute the yearly factor $1 + g = (y_T/y_0)^{1/T}$, and the growth rate $g$;
\item project: $y_{T+s} = y_T\,(1+g)^s$.
\end{enumerate}
\end{method}

\begin{example}[The United States, fitted on 1870 and 1929]\label{ex:g10:economic-models:us}
% figdata: trend_points (USA 4803, 11954, 58488)
Income per person in the United States was \usd{4803} in 1870 and \usd{11954} in 1929, 59 years
later (2011 dollars).
\begin{enumerate}
\item The ratio is $11954 / 4803 \approx 2.489$.
\item The yearly factor is $2.489^{1/59} \approx 1.0156$: a growth rate of about 1.56\,\% a year.
\item From 1929 to 2022 there are 93 years, so the forecast for 2022 is $11954 \times 1.0156^{93}
  \approx 11954 \times 4.219 \approx \usd{50434}$.
\end{enumerate}
The value observed in 2022 was \usd{58488}, so the observation was $58488 / 50434 \approx 1.16$
times the forecast: 16\,\% above it. The computer's line on the opening chart, which uses all sixty
points of 1870--1929 rather than two, gives 17\,\%.
% figdata: trend_fit (USA pct 117)
\end{example}

\begin{measure}[How a trend line is drawn on a times-ten scale]
A \emph{trend line} sums up a stretch of years by one straight line on a times-ten scale, that is,
by one constant growth rate. By hand, it is drawn through the first and last points of the
stretch, as in \cref{met:g10:economic-models:two-points}. That uses two numbers out of sixty, and
if either year was unusual (a war, a bad harvest) the line tilts with it. A computer instead
chooses the straight line that passes as close as possible to all the points at once; on the
opening chart, for the United States, it gives 1.53\,\% a year instead of the two-point 1.56\,\%, and for Canada 2.01\,\% instead of 1.87\,\%, enough to turn a near-perfect forecast into a miss of 13\,\%.
Two rules matter more than the method. First, say which years the line was fitted on. Second, judge
it only on other years.
% figdata: trend_fit (USA rate 1.53); trend_misses (CAN 2.01, pct 87); trend_points (CAN two-point 1.87)
\end{measure}

The model of the clerk's ruler can now be put in its box.

\begin{model}{Constant growth of income per person}
\assumes Income per person grows at the same yearly rate $g$ in every year; the rate of a past
  stretch of years carries on into the future.
\predicts On a times-ten scale, income per person follows a straight line; from the values of
  1870 and 1929, it gives a value for every later year, for example \usd{50434} for the United
  States in 2022.
\tested On 21 countries with long Maddison series, a line fitted on 1870--1929 lands within 25\,\%
  of the value observed in 2022 for four of them: Canada, Chile, the United States and
  Switzerland.
\breaks Through wars and depressions (the United States in the 1930s), in long stagnations
  (Argentina after 1930), and in fast catch-up, when a country grows much faster than before
  (Western Europe after 1950: most countries ended far above their line).
\end{model}
% figdata: trend_misses (CAN 87, CHL 88, USA 117, CHE 122; 21 countries)

\begin{omchart}{What the line of 1870--1929 missed: income per person observed in 2022, as a
percentage of the value the constant-growth line fitted on 1870--1929 predicted for 2022. At 100\,\%
the prediction was exact. Countries with at least 50 of the 60 years of 1870--1929 in the
database.}{Maddison Project Database 2023; GDP per person in 2011 dollars at purchasing-power parity}
\begin{tikzpicture}
\begin{axis}[omrecord, xbar, width=0.82\linewidth, height=10.5cm, xmin=0, xmax=1050,
  ymin=0.3, ymax=21.7, y dir=reverse, ytick=data, grid=none, xmajorgrids,
  yticklabels from table={figdata/grade-10/01-economic-models/out/trend_misses.csv}{name},
  xlabel={observed in 2022, as \% of the line's prediction}, bar width=2.6mm,
  y tick label style={font=\footnotesize}, xtick={0,100,200,...,1000},
  nodes near coords, nodes near coords style={font=\scriptsize}, point meta=x]
\addplot[fill=omDef!55, draw=omDef!80!black] table[x=pct, y=pos] {figdata/grade-10/01-economic-models/out/trend_misses.csv};
\draw[omThm!85!black, thick] (axis cs:100,0.3) -- (axis cs:100,21.7);
\end{axis}
\end{tikzpicture}
\end{omchart}

The chart of misses is as instructive as the opening chart. The four countries near 100\,\% did
not grow at the same rate all the time: the United States lost a third of its income per person
between 1929 and 1933 ($8048 / 11954 \approx 0.67$). They came back near their line. Most other countries did not: sixteen of
the twenty-one ended at more than one and a half times the value their line predicted, and one,
Argentina, at about half of it.
% figdata: obs_usa (1929 11954, 1933 8048: a fall of 33 %); trend_misses (16 at >= 153 %; ARG 56)
A fall of a third in four years, then a return to the line, is exactly what the constant-growth
model cannot describe and does not try to.

\begin{record}[A forecast that held, and one that did not]
% figdata: trend_points (ARG 6961 and USA 11954 in 1929; 18292 and 58488 in 2022)
In 1929, income per person in Argentina was \usd{6961} and in the United States \usd{11954}:
Argentina had 58\,\% of the American level. Both countries had grown at a little more than
1.5\,\% a year since 1870, and the constant-growth model predicted that the two would stay in
roughly the same proportion. In 2022, Argentina's income per person was \usd{18292} and the
American one \usd{58488}: Argentina had 31\,\% of the American level. The model held for one
country and broke for the other. What the record shows is \emph{that} it broke and \emph{when}:
the Argentine curve leaves its line in the 1930s and never comes back. It does not show
\emph{why}. A model that extends a line contains no cause at all, and so it cannot tell us what
changed in Argentina; that question needs a different kind of model, about institutions,
investment and policy, met in grade 12.
\end{record}

The constant-growth model gives a useful property that later chapters use often, when two
countries are compared.

\begin{proposition}[The ratio of two constant-growth paths]\label{prop:g10:economic-models:ratio}
If income per person grows at the constant rate $g_1$ in one country and $g_2$ in another, the
ratio of the first income to the second is multiplied by $(1+g_1)/(1+g_2)$ every year. It is
constant when $g_1 = g_2$, rises when $g_1 > g_2$ and falls when $g_1 < g_2$.
\end{proposition}
\begin{proof}
After $t$ years the two incomes are $a(1+g_1)^t$ and $b(1+g_2)^t$, where $a$ and $b$ are the
starting values. Their ratio is
\[
  \frac{a(1+g_1)^t}{b(1+g_2)^t} = \frac{a}{b}\left(\frac{1+g_1}{1+g_2}\right)^{t},
\]
so each extra year multiplies it by $(1+g_1)/(1+g_2)$. This factor equals 1 when $g_1 = g_2$, is
greater than 1 when $g_1 > g_2$, and smaller than 1 when $g_1 < g_2$.
\end{proof}

With the lines of the opening chart, the United States at 1.53\,\% and Argentina at 1.64\,\%, the
factor is $1.0153/1.0164 \approx 0.9989$: the model predicted that Argentina would very slowly
gain on the United States. It lost ground instead, which is what makes the Argentine line a clear
failure rather than a near miss.
% figdata: trend_fit (rates 1.53 and 1.64)

\section{Facts and values}\label{sec:g10:economic-models:facts-values}

``Argentina's income per person grew more slowly than Canada's after 1929'' and ``Argentina ought
to have grown faster'' are both sentences about Argentina, but they are of different kinds. The
first can be checked against the record; the second cannot, because it says what would be good,
and the record only says what happened.

\begin{definition}[Positive and normative statements]\label{def:g10:economic-models:positive}
A \emph{positive statement}\index{positive statement} says what is, was or would be; observation
can, at least in principle, show it true or false. A \emph{normative statement}\index{normative
statement} says what ought to be, or what is better or worse; it rests on values, and no
observation alone can settle it.
\end{definition}
\phantomsection\label{def:g10:economic-models:normative}

Three remarks keep the distinction useful.
\begin{itemize}
\item A \omterm{def:g10:economic-models:positive}{positive statement} can be false. ``Raising the price of a bus ticket by 10\,\% would halve
  the number of passengers'' is positive, and probably wrong; its kind is fixed by the fact that
  data could decide it.
\item A \omterm{def:g10:economic-models:positive}{positive statement} about a policy is still positive. ``A tax of one dalen per litre on
  petrol would reduce the petrol sold in Brenland by about 3\,\%'' is a prediction. Whether the tax
  is worth that reduction is normative.
\item Words can carry a verdict without saying so. ``Prices are too high'', ``a fair wage'' and
  ``an excessive profit'' each contain a normative judgment folded into what looks like a
  description. Written without the verdict, they become checkable: ``prices rose by 8\,\% while
  wages rose by 3\,\%''.
\end{itemize}

\begin{remark}[What this book does with values]\label{rem:g10:economic-models:values}
This book states models and tests them: its sentences are positive. When a choice depends on
values --- how much a gain to one group is worth against a loss to another --- the book says what
each option does, to whom and by how much, and leaves the weighing to the reader. Positive
economics does not make such choices for anyone; it makes them better informed.
\end{remark}

\section{Correlation and cause}\label{sec:g10:economic-models:cause}

The chart below puts a fact about farming beside a question. Countries whose farmers spread more
fertiliser on each hectare of land harvest more grain from each hectare. In the sixteen countries
that use less than \qty{10}{kg} of fertiliser per hectare, the middle (median) country harvests
about \qty{920}{kg} of cereals per hectare; in the fifty-six that use \qty{100}{kg} or more, about
\qty{4900}{kg}, more than five times as much.
% figdata: yield_groups (n 16, 57, 56; median 920, 2250, 4900)

\begin{omchart}{Cereal yield against fertiliser use, 129 countries with at least 100\,000 hectares
under cereals, averages of 2018--2020, both on times-ten scales. Fertiliser is measured in kilograms
of nutrients per hectare of arable land.}{World Bank, World Development Indicators (from FAO): cereal yield and fertiliser consumption}
\begin{tikzpicture}
\begin{loglogaxis}[omscatter, log basis x=10, log basis y=10, log ticks with fixed point,
  xtickten={0,...,3}, ytickten={2,...,5}, width=0.82\linewidth, height=7cm, xmin=1, xmax=1000,
  ymin=300, ymax=20000, xlabel={fertiliser (kg per hectare)}, ylabel={cereal yield (kg per hectare)}]
\addplot[omDef, only marks, mark=*, mark size=1.3pt, fill opacity=0.7]
  table[x=fert, y=yld] {figdata/grade-10/01-economic-models/out/yield_fert.csv};
\end{loglogaxis}
\end{tikzpicture}
\end{omchart}

The points rise from left to right: the two quantities are correlated. Does that mean that the
fertiliser \emph{causes} the higher yield, and by how much?

\begin{definition}[Causal effect]\label{def:g10:economic-models:causal-effect}
The \emph{causal effect}\index{causal effect} of a change in one variable on another is the change
in the second variable that the first change produces, \omterm{def:g10:economic-models:other-things-equal}{other things equal}.
\end{definition}

A correlation is what we see; a \omterm{def:g10:economic-models:causal-effect}{causal effect} is what we want to know, because only a \omterm{def:g10:economic-models:causal-effect}{causal effect}
says what would happen if someone changed the first variable. The scatter plot cannot give it, for
two reasons.

The first reason is a third variable. Countries that use a lot of fertiliser are, on the whole,
richer: their farmers also have tractors, irrigation, better seed, roads to the market and the
advice of agronomists. Each of these raises yields by itself. Part of the rise of the cloud of
points --- perhaps most of it, perhaps little --- comes from these other differences, which move
together with fertiliser.

\begin{definition}[Confounding factor]\label{def:g10:economic-models:confounder}
A \emph{confounding factor}\index{confounding factor} of the link between two variables is a third
variable that affects both of them. It creates a correlation between them even where neither causes
the other, and it makes a real \omterm{def:g10:economic-models:causal-effect}{causal effect} look larger or smaller than it is.
\end{definition}

The second reason is that the cause may run the other way. Where the soil is deep and the rain
reliable, a sack of fertiliser pays for itself several times over, so farmers buy more of it. High
yields then cause high fertiliser use, not the reverse.

\begin{definition}[Reverse causation]\label{def:g10:economic-models:reverse}
There is \emph{reverse causation}\index{reverse causation} when the variable thought of as the
effect causes, at least in part, the variable thought of as the cause.
\end{definition}

\begin{omfigure}
\begin{tikzpicture}[econ, v/.style={draw, rounded corners=1mm, minimum width=20mm, minimum
  height=7mm, font=\small, fill=white}, c/.style={->, very thick, omDef!80!black},
  k/.style={<->, thick, dashed, black!55}]
\begin{scope}
\node[v] (f1) at (0,0) {fertiliser}; \node[v] (y1) at (3.2,0) {yield};
\draw[c] (f1) -- (y1);
\node[font=\small\bfseries] at (1.6,-1.2) {(a) a cause};
\end{scope}
\begin{scope}[xshift=6.4cm]
\node[v] (z2) at (1.6,1.5) {farm wealth}; \node[v] (f2) at (0,0) {fertiliser}; \node[v] (y2) at (3.2,0) {yield};
\draw[c] (z2) -- (f2); \draw[c] (z2) -- (y2); \draw[k] (f2) -- (y2);
\node[font=\small\bfseries] at (1.6,-1.2) {(b) a confounding factor};
\end{scope}
\begin{scope}[yshift=-3.0cm, xshift=1.6cm]
\node[v] (f3) at (0,0) {fertiliser}; \node[v] (y3) at (3.2,0) {yield}; \node[v] (s3) at (6.4,0) {good soil};
\draw[c] (y3) -- (f3); \draw[c] (s3) -- (y3);
\node[font=\small\bfseries] at (3.2,-1.2) {(c) reverse causation};
\end{scope}
\end{tikzpicture}

{\small Three ways to get the same correlation. A solid arrow is a \omterm{def:g10:economic-models:causal-effect}{causal effect}; the dashed line in
(b) is a correlation without a direct cause. The scatter plot cannot tell the three apart; in
practice all three may be at work at once.}
\end{omfigure}

How, then, can a \omterm{def:g10:economic-models:causal-effect}{causal effect} be measured? The answer of agronomists is older than economics: take
a field, divide it into plots, and let chance decide which plots receive fertiliser.

\begin{definition}[Randomised experiment]\label{def:g10:economic-models:randomised-experiment}
A \emph{randomised experiment}\index{randomised experiment} is a comparison in which chance --- a
lottery, a coin, a table of random numbers --- decides which units (plots, people, firms, schools)
receive a treatment and which do not. The units that receive it form the \emph{treated group}, the
others the \emph{control group}.
\end{definition}

\begin{proposition}[What a lottery buys]\label{prop:g10:economic-models:lottery}
In a \omterm{def:g10:economic-models:randomised-experiment}{randomised experiment} with many units, the treated and control groups differ, on average, only
by the treatment. The difference between the average outcome of the treated group and the average
outcome of the control group then estimates the average \omterm{def:g10:economic-models:causal-effect}{causal effect} of the treatment, and the
estimate gets closer to it as the number of units grows.
\end{proposition}
\admitted

The proposition is proved in the university volumes, with the tools of probability. The idea can be
seen without them. When a coin decides, soil quality, the farmer's skill and the slope of the land
cannot influence who is treated, so they end up spread over the two groups in roughly the same way,
and the more plots there are, the closer the two groups are. Only the fertiliser is put on one group
and not the other. A \omterm{def:g10:economic-models:confounder}{confounding factor} has nothing to work with, and \omterm{def:g10:economic-models:reverse}{reverse causation} is ruled out:
the yield cannot decide who gets the fertiliser, because the coin did. Field trials of this kind
exist for fertiliser and many other farming practices, and a grade 11 chapter uses one.

\omterm{def:g10:economic-models:randomised-experiment}{Randomised experiments} are not always possible: no one can draw lots to decide which countries
get secure property rights. The grades 6--9 volume met the \emph{natural experiment}, in which
history or nature, not the researcher, splits a population in two, as the division of Korea did.
A natural experiment is only as good as one assumption: that, without the treatment, the two parts
would have moved in the same way. That assumption cannot be checked directly, since only one of the
two paths is ever seen for each part. What can be checked is the period before the split: if the
two parts were already moving apart, the assumption is in doubt.

\begin{remark}[Why the constant-growth model escapes all this]\label{rem:g10:economic-models:no-cause}
The model of the opening chart predicts without claiming any cause: it says ``if the past rate
carries on, then\ldots'', and nothing about why the rate is what it is. That makes it easy to test
and impossible to use for a decision, because it cannot say what would change if a country did
something differently. The models of the coming chapters are causal: they say what happens to
quantities when a price, a tax or a technology changes, \omterm{def:g10:economic-models:other-things-equal}{other things equal}. They are harder to
test, and the four-point frame is how this book keeps them honest.
\end{remark}

\section{Exercises}

\begin{exercise}[$\star$]\label{exo:g10:economic-models:1}
Read the opening chart. About how large was income per person in the United States and in Argentina
in 1929, and in 2022? In which country did the observed curve stay closer to its line?
\end{exercise}

\begin{exercise}[$\star$]\label{exo:g10:economic-models:2}
Say whether each statement is positive or normative.
\begin{enumerate}
\item In 2022, income per person in Canada was higher than in Chile.
\item Brenland ought to spend more on its railways.
\item If bus fares were free, the number of bus journeys would double.
\item A minimum wage is unfair to employers.
\item Income per person in Argentina grew more slowly after 1929 than before.
\end{enumerate}
\end{exercise}

\begin{exercise}[$\star$]\label{exo:g10:economic-models:3}
In the model $y_t = y_0\,(1+g)^t$, which variables are endogenous and which are exogenous? In the
bakery model $R = p\,q$ of \cref{ex:g10:economic-models:bakery}, which variable is endogenous?
\end{exercise}

\begin{exercise}[$\star$]\label{exo:g10:economic-models:4}
Read the chart of misses. How many countries ended between 75\,\% and 125\,\% of their line's
prediction? Which country ended furthest below its line, and which furthest above?
\end{exercise}

\begin{exercise}[$\star$]\label{exo:g10:economic-models:5}
``When the price of ice cream rises, fewer ice creams are sold.'' In July, the price of ice cream
rose in Brenland and more ice creams were sold than in May. Does this contradict the statement?
Use the phrase \emph{\omterm{def:g10:economic-models:other-things-equal}{other things equal}} in your answer.
\end{exercise}

\begin{exercise}[$\star\star$]\label{exo:g10:economic-models:6}
A country's income per person is \usd{20000} and grows at a constant 2\,\% a year. Compute it after
10 years and after 35 years. Roughly how many years does it take to double?
\end{exercise}

\begin{exercise}[$\star\star$]\label{exo:g10:economic-models:7}
% figdata: trend_points (CAN 2702, 8074, 45530)
Income per person in Canada was \usd{2702} in 1870 and \usd{8074} in 1929. Use
\cref{met:g10:economic-models:two-points} to find the growth rate of 1870--1929, rounded to two
decimal places in per cent, and the forecast for 2022. The value observed in 2022 was \usd{45530}.
By what percentage did the forecast miss?
\end{exercise}

\begin{exercise}[$\star\star$]\label{exo:g10:economic-models:8}
Two countries start with the same income per person. One grows at 2\,\% a year, the other at 1\,\%.
Use \cref{prop:g10:economic-models:ratio} to predict the ratio of their incomes after 70 years.
On the opening chart the lines grow at 1.53\,\% (United States) and 1.64\,\% (Argentina) a year.
By what factor does the model multiply the ratio of Argentine to American income over the 93 years
from 1929 to 2022? In 1929 that ratio was 0.582: what does the model predict for 2022?
\end{exercise}

\begin{exercise}[$\star\star$]\label{exo:g10:economic-models:9}
Read the scatter plot of fertiliser and yield. How many times larger is the median yield of the
countries using \qty{100}{kg} or more of fertiliser per hectare than that of the countries using
less than \qty{10}{kg}? Explain why this number is not the \omterm{def:g10:economic-models:causal-effect}{causal effect} of fertiliser on yield.
\end{exercise}

\begin{exercise}[$\star\star$]\label{exo:g10:economic-models:10}
Across countries, the number of mobile phones per hundred people is positively correlated with life
expectancy. Name a \omterm{def:g10:economic-models:confounder}{confounding factor}, and explain in two sentences how it creates the correlation.
\end{exercise}

\begin{exercise}[$\star\star$]\label{exo:g10:economic-models:11}
Firms that spend more on advertising have larger sales. Explain how causation could run, at least in
part, from sales to advertising. Then name one observation that would point to an effect of
advertising on sales.
\end{exercise}

\begin{exercise}[$\star\star\star$]\label{exo:g10:economic-models:12}
Give two observations, each of a kind that could be found in the data of a single country, that
would contradict the constant-growth model of income per person. For each, say which assumption of
the model it contradicts.
\end{exercise}

\begin{exercise}[$\star\star\star$]\label{exo:g10:economic-models:13}
A farming office in Brenland wants to know how much \qty{50}{kg} of fertiliser per hectare raises
the maize yield. It can work with 200 farms. Describe a \omterm{def:g10:economic-models:randomised-experiment}{randomised experiment}: how the farms are
chosen for the treatment, what is measured, and what number estimates the effect. Explain why
letting farmers volunteer for the fertiliser would give a misleading answer, and in which direction
it would most likely err.
\end{exercise}

\begin{exercise}[$\star\star\star$]\label{exo:g10:economic-models:14}
% figdata: trend_fit (pred2022 50167 and 32520; act2022 58488 and 18292); trend_points (1929)
The computer's lines of the opening chart predicted \usd{50167} for the United States and
\usd{32520} for Argentina in 2022; the observed values were \usd{58488} and \usd{18292}. In 1929
the values were \usd{11954} and \usd{6961}. Compute the ratio of American to Argentine income in
1929, the ratio the lines predicted for 2022, and the ratio observed in 2022. Which of the two
countries' forecasts contributed more to the error in the ratio? Explain why this observation
contradicts the model even though both countries grew.
\end{exercise}

\begin{exercise}[$\star\star\star$]\label{exo:g10:economic-models:15}
In 2019 a flood destroyed the harvest in the eastern half of a valley in Brenland, but not in the
western half. A researcher compares the incomes of eastern and western farm families in 2020 to
measure the effect of the flood. What must the researcher assume for this comparison to measure the
\omterm{def:g10:economic-models:causal-effect}{causal effect} of the flood? Which data from before 2019 would make the assumption more believable,
and which would cast doubt on it?
\end{exercise}

\section{Problem: Brenland's Trend}

\begin{problem}[{Weekend problem --- a forecast made in 2007, judged in 2025}]
\label{pb:g10:economic-models:1}
% figdata: brenland_gdppc (2000 24675, 2007 27027, 2009 25840, 2010 26170, 2014 27905,
% 2019 29997, 2020 28371, 2025 31470; Brenland is invented, code/shared/country)
At the end of 2007, an economist in Brenland is asked to forecast real GDP per person for 2025. The
statistics office publishes the series in dalens of 2015 per person. Some of its values, including
those published later, are:

\begin{center}
\small
\begin{tabular}{lrrrrrrrr}
\toprule
year & 2000 & 2007 & 2009 & 2010 & 2014 & 2019 & 2020 & 2025 \\
\midrule
dalens per person & 24\,675 & 27\,027 & 25\,840 & 26\,170 & 27\,905 & 29\,997 & 28\,371 & 31\,470 \\
\bottomrule
\end{tabular}
\end{center}

\noindent The economist uses the constant-growth model, fitted on 2000 and 2007 by
\cref{met:g10:economic-models:two-points}. Round growth rates to two decimal places in per cent and
yearly factors to four decimal places.

\medskip
\noindent\textbf{Part I --- The data.}
\begin{enumerate}
\item By what percentage did real GDP per person rise from 2000 to 2007?
\item By what percentage did it change from 2007 to 2009?
\item By what percentage did it change from 2019 to 2020?
\item In the economist's model, is real GDP per person of 2025 an endogenous or an \omterm{def:g10:economic-models:endogenous}{exogenous
  variable}? And the growth rate?
\item ``Brenland grew too slowly after 2008.'' Is this statement positive or normative? Rewrite it
  as a \omterm{def:g10:economic-models:positive}{positive statement} that the table can check.
\end{enumerate}

\medskip
\noindent\textbf{Part II --- The fitted rate.}
\begin{enumerate}[resume]
\item Compute the ratio of 2007 to 2000 and the yearly growth rate $g$ of 2000--2007.
\item Check the fit: starting from the value of 2000, what does the model give for 2007?
\item At this rate, roughly how many years does real GDP per person take to double?
\item What does the model assume about the years after 2007? Name one event that, if it happened,
  would break that assumption.
\end{enumerate}

\medskip
\noindent\textbf{Part III --- The projection.}
\begin{enumerate}[resume]
\item What does the model predict for 2009? By what percentage was the observed value of 2009 below
  this prediction?
\item What does the model predict for 2025?
\item Compute the forecast error in 2025: the observed value minus the forecast, as a percentage of
  the forecast.
\item The gap between forecast and observation in 2025 equals how many years of growth at the rate
  $g$, roughly? Compare $(1+g)^6$ with the ratio of the forecast to the observed value.
\end{enumerate}

\medskip
\noindent\textbf{Part IV --- What broke.}
\begin{enumerate}[resume]
\item Compute the growth rate of 2010--2019 from the two endpoints. Compare it with $g$.
\item Suppose the economist had restarted the forecast from the observed value of 2009, keeping the
  rate $g$. What would the forecast for 2025 have been, and what would the error have been?
\item Compare the errors of questions 12 and 15. Which assumption of the model failed more: that the
  rate stays the same, or that nothing interrupts the path?
\item The fall of 2020 was followed by a recovery. Does one year of falling income, by itself,
  contradict the model? What pattern over several years would?
\item A second economist forecasts in 2014, fitting the model on 2000 and 2014. Compute the rate and
  the forecast for 2025, and its error.
\item The 2014 forecast missed by less than the 2007 one. Does that show that its method was better?
  Explain.
\item State the result of the problem in one sentence: the forecast error of the 2007 forecast in
  2025, in per cent, and the event to which most of it is due.
\end{enumerate}
\end{problem}
